\chapter{Cauchy deconvolution from $\mathfrak U_1$ to $\mathfrak U_{1/2}$}\label{chap:cf}\chaptermark{Cauchy deconvolution to the line}

Under the Riemann hypothesis the level one density of Chapters~\ref{chap:lo} to~\ref{chap:cd} is the Cauchy smoothing of the zero counting measure, $u_1=P_{1/2}*\mathfrak U_{1/2}$ (equation~\eqref{es:eq:forward}), and the Thorin measure $\mathfrak U_\theta$ at every level $\theta\in(\frac12,1)$ is the same smoothing at scale $\theta-\frac12$. Can the deconvolution from $\mathfrak U_1$, which the Euler product makes accessible, to $\mathfrak U_{1/2}$, which is the object of the Riemann hypothesis, be carried out, and can the particle filtering methods of the book on Hilbert's sixth problem \cite{PSH626}, accelerated by generative Bayesian computation (GBC) \cite{PS23}, make progress on it? This chapter answers in three layers. The identities and criteria that relate the two measures are recalled with their labels (Section~\ref{cf:sec:recap}). A computational method follows: a particle filter whose time is the level, running the Cauchy flow backwards from $\theta=1$ towards $\theta=\frac12$ in a window of heights, and a GBC screen that replaces most of its work by one pass of a network. And the gap between local computation and the Riemann hypothesis, which is a statement for all heights, is recorded together with what a proof along these lines would still need (Section~\ref{cf:sec:gap}).

The chapter develops a note by Nick Polson \cite{P26cf}, with the computations done for this book. Several parts of the note are already in the book, in the same or in a stronger form, and are cited by label in Section~\ref{cf:sec:recap} and not restated: the density and the Hadamard sum, the transport identity, the deconvolution criterion (Theorem~\ref{tf:thm:deconv}, Proposition~\ref{cd:prop:critical}) and Thorin positivity as a zero free half plane (Theorem~\ref{cx:thm:positivity}, Lagarias \cite{Lagarias}, Sondow and Dumitrescu \cite{SD10}). The numbers of the chapter are of two kinds. The zeros and counts of the Davenport and Heilbronn function, the contour moments, the values of the density and the negative masses, the information numbers, the linear programs and polynomial tables and the nested sampling scores were recomputed for this book by scripts written from the mathematics (\texttt{code/cf\_dhzeros.py}, \texttt{cf\_moments.py}, \texttt{cf\_xi.py}, \texttt{cf\_info.py}, \texttt{cf\_lp.py}, \texttt{cf\_arith.py} and \texttt{cf\_ns.py}, with the data files \texttt{data/cf\_*.json}), and where the recomputation differed from the note the recomputed value is printed. The runs of the particle filter, of the GBC networks, of the concentration experiments and of the nested sampler are runs of the note's own code, which was not run for the book: they are printed as the note reports them and are marked computed by the author.

Status words follow the rest of the book: proved, computed, conditional, heuristic or open. Statements at the level $\frac78$ use the quasi Riemann hypothesis $\mathrm{QRH}(\frac78)$, announced by OpenAI \cite{OAI26a} and formalised in Lean, and not refereed (Section~\ref{lv:sec:qrh}); every one of them is conditional on it. Nothing in this chapter proves the Riemann hypothesis.

\begin{enumerate}
\item The notation and the cited criteria (Section~\ref{cf:sec:recap}): the right limit at a level through a zero (Remark~\ref{cf:rem:limit}), the transport identity (Proposition~\ref{cf:prop:transport}, proved), and local positivity of the density from any zero free region that narrows with the height, in particular the regions of de la Vall\'ee Poussin and of Vinogradov and Korobov (Corollary~\ref{cf:cor:localpos}, proved). The formula of Balazard, Saias and Yor is proved at every level in Chapter~\ref{chap:bsy}.
\item Two test functions (Section~\ref{cf:sec:test}): $\zeta$ and the Davenport and Heilbronn function, the control with zeros off the line, whose four zeros off the line below height $200$ are computed and certified (Proposition~\ref{cf:prop:ends}, proved; the zeros computed and certified in ball arithmetic).
\item A particle filter along the Cauchy flow (Section~\ref{cf:sec:pf}): the model and the sampler, the detectability of a symmetric pair (Proposition~\ref{cf:prop:detect}, proved), and the information carried by each level, which is first order for the filter as run and second order for a pair (heuristic).
\item GBC (Section~\ref{cf:sec:gbc}): the contour moments, which determine the zeros in a rectangle exactly (Proposition~\ref{cf:prop:moments}, proved; Table~\ref{cf:tab:dl} computed), the learned screen and its calibration (computed by the author), and its use to initialise the filter.
\item Concentration of the particle approximation and the window prediction of $u_{3/4}$ (Sections~\ref{cf:sec:conc} and~\ref{cf:sec:u34}): a concentration inequality for the idealised system (Theorem~\ref{cf:thm:conc}, proved), which does not cover the samplers run, and the prediction from a fit to two levels (computed by the author, Proposition~\ref{cf:prop:determine}, proved).
\item Negative mass and arithmetic inequalities (Sections~\ref{cf:sec:negmass} and~\ref{cf:sec:arith}): the unconditional bound on the negative part of $\mathfrak U_\theta$ in terms of zero density (Theorem~\ref{cf:thm:negmass}, proved), the converse results, the zero free regions that follow from positivity of the density at the heights $0$, $t$ and $2t$ (Theorem~\ref{cf:thm:dvp}, proved), the barrier that the Davenport and Heilbronn function sets (Proposition~\ref{cf:prop:barrier}, proved), and the comparison of what multilevel schemes can give: a computed estimate of about $1/(11.3\log\gamma)$, a proved cap, and the constants that pair each zero with its reflection (Propositions~\ref{cf:prop:barrier2} and~\ref{cf:prop:stechkin}).
\item Nested sampling over heights and the open questions (Sections~\ref{cf:sec:ns} and~\ref{cf:sec:gap}): the normalised margin of $\zeta$ and the zeros of the Davenport and Heilbronn function near height $1000$ (computed), and what would close the gap.
\end{enumerate}

\section{What the book already proves}\label{cf:sec:recap}

\subsection{Notation}

The notation is the book's. The completed function is $\xi$ of Section~\ref{lo:sec:level}, the zeros are $\rho=\sigma_\rho+i\gamma_\rho$ with multiplicity, a level is $\theta$ (the real part of $s=\theta+iv$), and
\[
u_\theta(v)=\frac1\pi\Re\frac{\xi'}\xi(\theta+iv)=\frac1\pi\,\partial_\sigma\log\lvert\xi(\sigma+iv)\rvert\Big|_{\sigma=\theta}
\]
is the density of Theorem~\ref{lo:thm:density}. The measure $\mathfrak U_\theta=\sum_\rho\mathrm{Cauchy}(\gamma_\rho,\theta-\sigma_\rho)$ of equation~\eqref{lo:eq:Utheta} lives on the line, in the variable $v$, and carries no factor $\frac12$: $\mathfrak U_\theta=2\widetilde U_\theta$ with the $\widetilde U_\theta$ of Theorem~\ref{lv:thm:semigroup}, and the Thorin measure $U_\theta$ of Section~\ref{lv:sec:thorinform}, in the variable $x=v^2$, is the image of the restriction of $\mathfrak U_\theta$ to $v>0$. The symbol $U_\theta$ is used only for it. The Cauchy density is $P_a(x)=(a/\pi)/(a^2+x^2)$ for $a>0$, with $P_a*P_b=P_{a+b}$ and $\widehat{P_a}(\omega)=e^{-a\lvert\omega\rvert}$, and we write $Q_a(x)=(a/\pi)/(a^2+x^2)$ for every real $a\ne0$, so that $Q_a=-P_{\lvert a\rvert}$ is a Cauchy density with negative mass when $a<0$, and $Q_0(x)=0$ for $x\ne0$. For $\theta$ not equal to the real part of a zero at the height in question the Hadamard sum (equation~\eqref{es:eq:lorentz} at $\theta=1$, the proof of Theorem~\ref{cx:thm:positivity} at every level) reads
\begin{equation}\label{cf:eq:mg}
u_\theta(v)=\sum_\rho Q_{\theta-\sigma_\rho}(v-\gamma_\rho),
\end{equation}
absolutely convergent because $\sum_\rho\lvert\gamma_\rho\rvert^{-2}<\infty$. The function $(\sigma,v)\mapsto\pi^{-1}\Re(\xi'/\xi)(\sigma+iv)$ is harmonic off the zeros, as the real part of an analytic function, and positivity of $u_\theta$ at one level is the statement $\partial_\sigma\log\lvert\xi\rvert\ge0$ there; positivity at every level $\theta>\sigma_0$ is the monotone increase of $\lvert\xi(\sigma+iv)\rvert$ in $\sigma>\sigma_0$ studied by Sondow and Dumitrescu \cite{SD10} and by Matiyasevich, Saidak and Zvengrowski \cite{MSZ14}.

\begin{remark}[a level through a zero]\label{cf:rem:limit}
At a level $\theta=\sigma_\rho$ the term $Q_{\theta-\sigma_\rho}(v-\gamma_\rho)$ of~\eqref{cf:eq:mg} has no meaning at $v=\gamma_\rho$. For every $\theta>\frac12$ we take $\mathfrak U_\theta$ to be the vague limit of $u_{\theta'}(v)\,dv$ as $\theta'\downarrow\theta$. Each term converges vaguely, and the far terms are dominated by $\lvert\gamma_\rho\rvert^{-2}$ uniformly for $\theta'$ in a compact set, so the limit exists; it is the measure of equation~\eqref{lo:eq:Utheta} with an atom at the ordinate of every zero with $\sigma_\rho=\theta$ and a negative Cauchy term $Q_{\theta-\sigma_\rho}$ for every zero with $\sigma_\rho>\theta$. At $\theta=\frac12$ it is $\mathfrak U_{1/2}=\sum_{\sigma_\rho=1/2}\delta_{\gamma_\rho}$ plus the contribution of the zeros off the line, which vanishes (Proposition~\ref{cf:prop:ends}).
\end{remark}

\subsection{The transport identity and the criteria}

\begin{proposition}[Cauchy transport]\label{cf:prop:transport}
For $\theta>\frac12$ and $\theta+iv$ not a zero,
\begin{equation}\label{cf:eq:transport}
u_\theta=P_{\theta-1/2}*\mathfrak U_{1/2}+R_\theta,\qquad R_\theta(v)=\sum_{\sigma_\rho>1/2}\bigl[Q_{\theta-\sigma_\rho}+Q_{\theta-1+\sigma_\rho}\bigr](v-\gamma_\rho),
\end{equation}
the sum over the zeros off the critical line, each taken together with its reflection $1-\bar\rho$, which has the same ordinate. In particular $u_1=P_{1/2}*\mathfrak U_{1/2}+R_1$: the density at level one is the Cauchy smoothing of $\mathfrak U_{1/2}$ at scale $\frac12$ plus one Cauchy pair $P_{1-\sigma_\rho}+P_{\sigma_\rho}$ for each such zero. The pair is positive at $\theta=1$ and is a signed term only for $\theta<\sigma_\rho$.
\end{proposition}

\begin{proof}
Split the sum~\eqref{cf:eq:mg} into the zeros on the line, each contributing $P_{\theta-1/2}(\cdot-\gamma_\rho)=P_{\theta-1/2}*\delta_{\gamma_\rho}$, and the pairs $\rho$, $1-\bar\rho$ off it.
\end{proof}

This is Theorem~\ref{lv:thm:semigroup} above the abscissa of the last zero, and equation~\eqref{es:eq:forward} under the Riemann hypothesis, where $R_\theta=0$; what is stated here is the identity of densities at every level, the second term being the defect. In Fourier terms a zero on the line contributes $e^{-(\theta-1/2)\lvert\omega\rvert}e^{i\gamma\omega}$, while a pair at distance $\delta=\sigma_\rho-\frac12$ contributes $2\cosh(\delta\omega)\,e^{-(\theta-1/2)\lvert\omega\rvert}e^{i\gamma\omega}$ for $\theta>\sigma_\rho$. Deconvolving by $P_{1/2}$ from $\theta=1$ multiplies by $e^{\lvert\omega\rvert/2}$, and the result $2\cosh(\delta\omega)$ is not the Fourier transform of any tempered distribution (Chapter~\ref{chap:cd}). This is the whole problem in one line.

Three statements are known and are not reproved. First, the deconvolution criterion: the Riemann hypothesis holds if and only if $u_1=P_{1/2}*\nu$ for a positive measure $\nu$ with $\int(1+v^2)^{-1}\nu(dv)<\infty$, and then $\nu=\mathfrak U_{1/2}$ (Theorem~\ref{tf:thm:deconv}(b); the critical width of a positive datum is Proposition~\ref{cd:prop:critical}, with the Herglotz representation of Theorem~\ref{cd:thm:herglotz}, in the context of Chapter~\ref{chap:cd}). Second, Thorin positivity is a zero free half plane: for $\frac12\le\sigma_0<1$ the measures $\mathfrak U_\tau$ are positive for every $\tau>\sigma_0$ if and only if $\xi$ has no zero with $\sigma_\rho>\sigma_0$, and for $\sigma_0=\frac12$ this is the theorem of Lagarias \cite{Lagarias} in the form of Sondow and Dumitrescu \cite{SD10} (Theorem~\ref{cx:thm:positivity}); positivity at the single level $\sigma_0$ is necessary and is not known to be sufficient (Remark~\ref{rem:ch2:oneline}). Third, at levels $\theta\ge1$ the density is positive unconditionally, $u_\theta=\sum_\rho P_{\theta-\sigma_\rho}(\cdot-\gamma_\rho)>0$ (Theorem~\ref{lo:thm:density}, Corollary~\ref{lv:cor:uncond}). The pairing that underlies every statement of this chapter about a fixed height is that of Proposition~\ref{prop:signature}(ii), used in Corollary~\ref{cor:dbeta:window}(1) and Lemma~\ref{cx:lem:pairing}.

\subsection{Local positivity from a zero free region}

A zero free region that narrows with the height is not a half plane, so Theorem~\ref{cx:thm:positivity} does not apply to it directly. The pairing gives the local statement.

\begin{corollary}[local positivity from a zero free region]\label{cf:cor:localpos}
Let $\eta:[0,\infty)\to(0,\frac12]$ be nonincreasing and suppose that every zero satisfies $\sigma_\rho\le1-\eta(\lvert\gamma_\rho\rvert)$. Then $u_\theta(v)>0$ whenever $\theta>\frac12$, $\theta\ge1-\eta(\lvert v\rvert+\frac12)$ and $\xi(\theta+iv)\ne0$.
\end{corollary}

\begin{proof}
Group the zeros in the Hadamard sum~\eqref{es:eq:lorentz} into orbits of $\rho\mapsto1-\bar\rho$. As in the proof of Corollary~\ref{cor:dbeta:window}(1) (Proposition~\ref{prop:signature}(ii)), a zero on the critical line contributes $P_{\theta-1/2}(v-\gamma_\rho)>0$, an orbit with $\sigma_\rho\le\theta$ contributes a nonnegative amount, and an orbit with $\sigma_\rho>\theta$ is negative only where $\lvert v-\gamma_\rho\rvert<\sqrt{(\sigma_\rho-\theta)(\sigma_\rho+\theta-1)}<\frac12$. A zero with $\sigma_\rho>\theta$ and $\lvert v-\gamma_\rho\rvert<\frac12$ would have $\lvert\gamma_\rho\rvert<\lvert v\rvert+\frac12$, hence $\sigma_\rho\le1-\eta(\lvert\gamma_\rho\rvert)\le1-\eta(\lvert v\rvert+\frac12)\le\theta$, a contradiction. So every orbit contributes a nonnegative amount, and the zero $\frac12+i\gamma_1$ a positive one.
\end{proof}

With the region of de la Vall\'ee Poussin ($\eta(T)$ a constant over $\log T$, Proposition~\ref{tx3:prop:dvp}) or of Vinogradov and Korobov ($\eta(T)=c(\log T)^{-2/3}(\log\log T)^{-1/3}$ \cite{Vinogradov58}), taken as $\frac12$ at small heights, this is positivity of $u_\theta$ at height $v$ for $\theta$ within that width of the line $\Re s=1$, the width taken at height $\lvert v\rvert+\frac12$. It is local in $v$, as it must be: the width shrinks with height, and positivity at the single level $\theta$ does not by itself give a zero free half plane. Positivity of $\mathfrak U_\tau$ for every $\tau>\sigma$ is $\mathrm{QRH}(\sigma)$ (Theorem~\ref{cx:thm:positivity}), while positivity at the single level $\sigma$ is necessary and not known to be sufficient (Remark~\ref{rem:ch2:oneline}).

\subsection{Two cited identities}

Two further identities of the note are in the book or in the next chapter. The logarithmic version of the transport, in which the defect of $\log\lvert h\rvert$ becomes a Green potential of the zeros off the line, is the formula of Balazard, Saias and Yor \cite{BSY99}: Chapter~\ref{chap:bsy} proves the formula of Balazard, Saias and Yor at every level $\theta$ (Theorem~\ref{bsy:thm:bsy}); here it enters only through Proposition~\ref{cf:prop:barrier}. And the Cauchy density is the harmonic measure of a half plane: $\mathrm{Cauchy}(\gamma,a)$ is the law of $\gamma+B(T_a)$ for a Brownian motion $B$ and the hitting time $T_a$ of level $a$ by an independent Brownian motion (the closing remark of Section~\ref{lo:sec:gbc}, Section~\ref{bp:sec:planar}), so that under the Riemann hypothesis $u_\theta(v)$ is the summed hitting density of the critical zeros for planar Brownian motion started at $\theta+iv$. The forward flow is a Markov semigroup; the backward flow has no positive kernel, it is an inverse problem, and the natural particle method for it is a Bayesian filter and not a forward simulation.
